\begin{definition}[Structure coefficients indexed by a BNT]%
    \label{def:mpdo_bnt_label_coefficient_family}
    \lean{MPOTensor.BNTLabelCoefficientFamily}
    \lean{MPOTensor.BNTLabelCoefficientFamily.ofChi}
    \lean{MPOTensor.BNTLabelCoefficientFamily.ofChi_coeff}
    \lean{MPOTensor.BNTLabelCoefficientFamily.ofChi_coeff_eq_trace_matrix_pow}
    \leanok
    Let $\Lambda$ be a fixed set of BNT labels. Consider complex coefficients
    $c^{(L)}_{\alpha,\beta,\gamma}$ for $L\in\N$ and
    $\alpha,\beta,\gamma\in\Lambda$. In
    \cite[Theorem~4.14(ii)]{Cirac2016MPDO_arXiv}, these are the structure
    coefficients for the products $O_L(M_\alpha)O_L(M_\beta)$ at the same
    length $L$, rather than for the blocked-basis multiplication
    $\mathcal A_n\times\mathcal A_n\to\mathcal A_{2n}$.
    A family of diagonal matrices $\chi_{\alpha,\beta,\gamma}$ determines
    coefficients
    $c^{(L)}_{\alpha,\beta,\gamma}
        = \tr(\chi_{\alpha,\beta,\gamma}^{L})
        = \sum_r\chi_{\alpha,\beta,\gamma,r}^{\,L}$.
    The positive diagonal matrices constructed in
    \cite[Appendix~C.4]{Cirac2016MPDO_arXiv} give coefficients of this form;
    the definition for a general diagonal family does not require positivity.
\end{definition}

\begin{definition}[Operators indexed by a BNT]%
    \label{def:mpdo_bnt_label_operator_family}
    \lean{MPOTensor.BNTLabelOperatorFamily}
    \leanok
    For each length $L$, let $O_L(M_\alpha)$ be operators indexed by
    the fixed BNT labels $\alpha$. Their ambient complex vector space,
    equipped with multiplication, may depend on $L$ and is not identified
    with a chosen blocked support algebra.
    % Formula and source passage: lines 962--967 define O_L(M) as the trace of
    % the L-fold product. In components this is the equality displayed below.
    % Ink-to-index after demotion: there is no diagram ink. The multi-indices
    % \mathbf a and \mathbf b are free horizontal indices; the matrix indices
    % of the factors \widehat M_alpha^{a_r b_r}=M_alpha^{(a_r,b_r)} are
    % contracted by the trace.
    % Contraction set: precisely the cyclic matrix-index contraction encoded by
    % \tr; no physical index is summed.
    % Type verdict and boundary signature: the result is an operator on the
    % length-L horizontal physical space, with components indexed by
    % (\mathbf a,\mathbf b). An exact displayed formula is preferable here
    % because the old vertical-closure picture obscured which legs stayed free.
    Write $\widehat M_\alpha^{ab}=M_\alpha^{(a,b)}$ for the vertical reading.
    The vertical indices of the $L$ copies are contracted cyclically, while
    the horizontal legs remain open, as in
    \cite[lines~962--967]{Cirac2016MPDO_arXiv}.
    For $\mathbf a=(a_1,\ldots,a_L)$ and $\mathbf b=(b_1,\ldots,b_L)$,
    \begin{align}
        [O_L(M_\alpha)]_{\mathbf a,\mathbf b}
         & =
        \tr\!\left(\widehat M_\alpha^{a_1b_1}\cdots
        \widehat M_\alpha^{a_Lb_L}\right).
        \notag
    \end{align}
\end{definition}

\begin{definition}[Product of BNT operators at the same length]%
    \label{def:mpdo_bnt_label_same_length_product_form}
    \lean{MPOTensor.BNTLabelOperatorFamily.HasSameLengthProductForm}
    \lean{MPOTensor.BNTLabelOperatorFamily.HasSameLengthProductForm.eq_sum}
    \leanok
    \uses{def:mpdo_bnt_label_operator_family, def:mpdo_bnt_label_coefficient_family}
    For a finite set of BNT labels, the operators $O_L(M_\alpha)$ satisfy
    the product identity with coefficients $c^{(L)}_{\alpha,\beta,\gamma}$
    when, for every positive length $L$ and all labels $\alpha,\beta$,
    \begin{align}
        O_L(M_\alpha)O_L(M_\beta)
         & =
        \sum_\gamma
        c^{(L)}_{\alpha,\beta,\gamma} O_L(M_\gamma).
        \notag
    \end{align}
    This is the product identity appearing in
    \cite[Theorem~4.14(ii)]{Cirac2016MPDO_arXiv}, stated independently of
    the blocked-basis multiplication
    $\mathcal A_n\times\mathcal A_n\to\mathcal A_{2n}$.
\end{definition}

\begin{definition}[Scalars indexed by a BNT]%
    \label{def:mpdo_bnt_label_trace_scalar_family}
    \lean{MPOTensor.BNTLabelTraceScalarFamily}
    \leanok
    Let $m_\alpha\in\C$ be scalars indexed by the fixed BNT labels
    $\alpha$. For the positive diagonal matrices $\mu_\alpha$ in the
    vertical decomposition, these scalars are $m_\alpha=\tr(\mu_\alpha)$,
    as in the idempotent condition of
    \cite[Theorem~4.14(ii)]{Cirac2016MPDO_arXiv}.
\end{definition}

\begin{definition}[Idempotent for the multiplication induced by the coefficients]%
    \label{def:mpdo_bnt_label_idempotent_coefficient_form}
    \lean{MPOTensor.BNTLabelTraceScalarFamily.HasIdempotentCoefficientForm}
    \lean{MPOTensor.BNTLabelTraceScalarFamily.HasIdempotentCoefficientForm.eq_sum}
    \leanok
    \uses{def:mpdo_bnt_label_trace_scalar_family, def:mpdo_bnt_label_coefficient_family}
    For a finite set of BNT labels, the vector $m=(m_\alpha)_\alpha$ is an
    idempotent for the multiplication induced by $c^{(1)}$ when, for every
    label $\gamma$,
    \begin{align}
        m_\gamma
         & =
        \sum_{\alpha,\beta}
        c^{(1)}_{\alpha,\beta,\gamma}m_\alpha m_\beta.
        \notag
    \end{align}
    This is the idempotent condition in
    \cite[Theorem~4.14(ii)]{Cirac2016MPDO_arXiv}, stated separately from
    the trace-power formula.
\end{definition}

\begin{definition}[Assignment of BNT labels to blocked bases]%
    \label{def:mpdo_bnt_blocked_basis_label_assignment}
    \lean{MPOTensor.BNTBlockedBasisLabelAssignment}
    \lean{MPOTensor.BNTBlockedBasisLabelAssignment.blockedLabel}
    \leanok
    \uses{def:mpdo_blocked_structure_coefficients}
    An assignment of BNT labels consists, for every positive blocking length
    $n$, of a map $\sigma_n$ from the chosen basis indices of
    $\mathcal A_n$ and a map $\tau_n$ from the chosen basis indices of
    $\mathcal A_{2n}$ to the fixed set of BNT labels. The combined
    map $\sigma_n \sqcup \tau_n$ is defined on the disjoint union of these
    two index sets.
\end{definition}

\begin{definition}[Comparison of BNT structure coefficients with blocked coefficients]%
    \label{def:mpdo_bnt_blocked_basis_coefficient_comparison}
    \lean{MPOTensor.BNTBlockedBasisCoefficientComparison}
    \lean{MPOTensor.BNTBlockedBasisCoefficientComparison.sourceLabel}
    \lean{MPOTensor.BNTBlockedBasisCoefficientComparison.targetLabel}
    \lean{MPOTensor.BNTBlockedBasisCoefficientComparison.blocked_coeff_eq}
    \leanok
    \uses{def:mpdo_bnt_label_coefficient_family, def:mpdo_blocked_structure_coefficients, def:mpdo_bnt_blocked_basis_label_assignment}
    A comparison between the BNT structure coefficients and the chosen
    blocked-basis coefficients consists of an assignment of BNT labels
    together with the equality
    \begin{align}
        c^{(n)}_{i,j,k}
         & =
        c^{(n)}_{\sigma_n(i),\sigma_n(j),\tau_n(k)}
        \notag
    \end{align}
    for every $n>0$ and all basis indices $i,j$ of $\mathcal A_n$ and $k$
    of $\mathcal A_{2n}$.
    On the left, $c^{(n)}_{i,j,k}$ is the coefficient of the $k$-th basis
    element of $\mathcal A_{2n}$ in the blocked product. This comparison is
    an additional hypothesis, not the product identity for BNT operators
    at the same length.
\end{definition}

\begin{definition}[Trace-power formula at positive lengths]%
    \label{def:mpdo_bnt_label_positive_length_chi_trace_power_form}
    \lean{MPOTensor.BNTLabelCoefficientFamily.HasPositiveLengthChiTracePowerForm}
    \leanok
    \uses{def:mpdo_bnt_label_coefficient_family, def:mpdo_diagonal_chi_family}
    The coefficients $c^{(L)}_{\alpha,\beta,\gamma}$ satisfy the trace-power
    formula for a family of diagonal matrices $\chi_{\alpha,\beta,\gamma}$
    when, for every positive length $L$ and all labels,
    \begin{align}
        c^{(L)}_{\alpha,\beta,\gamma}
         & =
        \tr(\chi_{\alpha,\beta,\gamma}^{L}).
        \notag
    \end{align}
    The same matrix $\chi_{\alpha,\beta,\gamma}$ is used for all positive
    $L$. Its diagonal entries may be complex; positivity is not required
    here.
\end{definition}

\begin{lemma}[Structure coefficient as a trace power at positive length]%
    \label{lem:mpdo_bnt_label_chi_eq_trace_matrix_pow}
    \lean{MPOTensor.BNTLabelCoefficientFamily.HasPositiveLengthChiTracePowerForm.eq_trace_matrix_pow}
    \leanok
    \uses{def:mpdo_bnt_label_positive_length_chi_trace_power_form, thm:mpdo_trace_matrix_pow}
    Under the trace-power condition at positive lengths, for every $L>0$,
    $c^{(L)}_{\alpha,\beta,\gamma}
        = \tr(\chi_{\alpha,\beta,\gamma}^{L})$.
\end{lemma}
\begin{proof}\leanok
    Combine the defining positive-length trace-power identity with the trace
    formula for powers of a diagonal matrix.
\end{proof}

\begin{lemma}[Coefficients defined by traces of diagonal matrix powers]%
    \label{lem:mpdo_bnt_label_of_chi_positive_length_trace_power}
    \lean{MPOTensor.BNTLabelCoefficientFamily.ofChi_hasPositiveLengthChiTracePowerForm}
    \leanok
    \uses{def:mpdo_bnt_label_coefficient_family, def:mpdo_bnt_label_positive_length_chi_trace_power_form}
    The coefficients defined by
    $c^{(L)}_{\alpha,\beta,\gamma}=\tr(\chi_{\alpha,\beta,\gamma}^{L})$
    for a family of diagonal matrices satisfy the trace-power condition at
    positive lengths for those matrices.
\end{lemma}
\begin{proof}\leanok
    This follows from the definition of the coefficients.
\end{proof}

\begin{definition}[Positive diagonal matrices giving the structure coefficients]%
    \label{def:mpdo_positive_bnt_label_chi_trace_power_form}
    \lean{MPOTensor.PositiveBNTLabelChiTracePowerForm}
    \lean{MPOTensor.PositiveBNTLabelChiTracePowerForm.ofChi}
    \leanok
    \uses{def:mpdo_bnt_label_positive_length_chi_trace_power_form, def:mpdo_diagonal_chi_family, lem:mpdo_bnt_label_of_chi_positive_length_trace_power}
    For given coefficients $c^{(L)}_{\alpha,\beta,\gamma}$, consider diagonal
    matrices $\chi_{\alpha,\beta,\gamma}$, independent of $L$, with strictly
    positive diagonal entries, such that
    $c^{(L)}_{\alpha,\beta,\gamma}=\tr(\chi_{\alpha,\beta,\gamma}^{L})$
    for every $L>0$ and all labels. This is the coefficient condition in
    \cite[Theorem~4.14(ii)]{Cirac2016MPDO_arXiv}.
    If the coefficients are defined by these traces, positivity of the
    diagonal entries suffices.
\end{definition}

\begin{lemma}[Trace powers of the positive diagonal matrices]%
    \label{lem:mpdo_positive_bnt_label_chi_eq_trace_matrix_pow}
    \lean{MPOTensor.PositiveBNTLabelChiTracePowerForm.eq_trace_pow}
    \leanok
    \uses{def:mpdo_positive_bnt_label_chi_trace_power_form, lem:mpdo_bnt_label_chi_eq_trace_matrix_pow}
    If the coefficients have a trace-power representation by positive
    diagonal matrices $\chi_{\alpha,\beta,\gamma}$ independent of $L$, then
    for every $L>0$,
    $c^{(L)}_{\alpha,\beta,\gamma}
        =\tr(\chi_{\alpha,\beta,\gamma}^{L})$.
\end{lemma}
\begin{proof}\leanok
    Apply the trace reformulation of the trace-power identity at positive
    lengths.
\end{proof}

\begin{theorem}[Product of BNT operators with trace-power coefficients]%
    \label{thm:mpdo_bnt_label_same_length_product_eq_sum_chi_trace_pow}
    \lean{MPOTensor.BNTLabelOperatorFamily.HasSameLengthProductForm.eq_sum_chi_trace_pow}
    \leanok
    \uses{def:mpdo_bnt_label_same_length_product_form, lem:mpdo_positive_bnt_label_chi_eq_trace_matrix_pow}
    Suppose the product identity for BNT operators at the same length holds
    and its coefficients have a trace-power representation by positive
    diagonal matrices $\chi_{\alpha,\beta,\gamma}$ independent of $L$.
    Then, for every $L>0$,
    \begin{align}
        O_L(M_\alpha)O_L(M_\beta)
         & =
        \sum_\gamma
        \tr(\chi_{\alpha,\beta,\gamma}^{\,L})
        O_L(M_\gamma).
        \notag
    \end{align}
\end{theorem}
\begin{proof}\leanok
    Substitute the positive-length trace-power formula into the same-length BNT
    product expansion.
\end{proof}

\begin{theorem}[Product of BNT operators for coefficients defined by traces]%
    \label{thm:mpdo_bnt_label_same_length_product_eq_sum_of_chi_trace_pow}
    \lean{MPOTensor.BNTLabelOperatorFamily.HasSameLengthProductForm.eq_sum_ofChi_trace_pow}
    \leanok
    \uses{def:mpdo_bnt_label_same_length_product_form, def:mpdo_bnt_label_coefficient_family}
    If the product identity for BNT operators at the same length holds with
    coefficients defined by traces of powers of diagonal matrices
    $\chi_{\alpha,\beta,\gamma}$, then for every $L>0$,
    \begin{align}
        O_L(M_\alpha)O_L(M_\beta)
         & =
        \sum_\gamma
        \tr(\chi_{\alpha,\beta,\gamma}^{\,L})
        O_L(M_\gamma).
        \notag
    \end{align}
\end{theorem}
\begin{proof}\leanok
    Substitute the defining trace-power formula for the coefficients into the
    same-length BNT product expansion.
\end{proof}

\begin{theorem}[Idempotent condition in terms of traces]%
    \label{thm:mpdo_bnt_label_idempotent_eq_sum_chi_trace}
    \lean{MPOTensor.BNTLabelTraceScalarFamily.HasIdempotentCoefficientForm.eq_sum_chi_trace}
    \leanok
    \uses{def:mpdo_bnt_label_idempotent_coefficient_form, lem:mpdo_positive_bnt_label_chi_eq_trace_matrix_pow}
    If $m$ is an idempotent for the multiplication induced by $c^{(1)}$ and
    the coefficients have a trace-power representation by positive diagonal
    matrices $\chi_{\alpha,\beta,\gamma}$ independent of $L$, then
    \begin{align}
        m_\gamma
         & =
        \sum_{\alpha,\beta}
        \tr(\chi_{\alpha,\beta,\gamma})
        m_\alpha m_\beta.
        \notag
    \end{align}
\end{theorem}
\begin{proof}\leanok
    Substitute the length-one trace-power formula into the idempotent
    coefficient identity.
\end{proof}

\begin{theorem}[Idempotent condition for coefficients defined by traces]%
    \label{thm:mpdo_bnt_label_idempotent_eq_sum_of_chi_trace}
    \lean{MPOTensor.BNTLabelTraceScalarFamily.HasIdempotentCoefficientForm.eq_sum_ofChi_trace}
    \leanok
    \uses{def:mpdo_bnt_label_idempotent_coefficient_form, def:mpdo_bnt_label_coefficient_family}
    If $m$ is an idempotent for the multiplication induced by $c^{(1)}$,
    where the coefficients are defined by traces of powers of diagonal
    matrices $\chi_{\alpha,\beta,\gamma}$, then
    \begin{align}
        m_\gamma
         & =
        \sum_{\alpha,\beta}
        \tr(\chi_{\alpha,\beta,\gamma})
        m_\alpha m_\beta.
        \notag
    \end{align}
\end{theorem}
\begin{proof}\leanok
    Substitute the defining length-one trace formula into the idempotent
    coefficient identity.
\end{proof}

\begin{lemma}[Blocked coefficients inherit BNT trace powers]%
    \label{lem:mpdo_bnt_blocked_basis_coefficient_eq_trace_pow}
    \lean{MPOTensor.BNTBlockedBasisCoefficientComparison.blocked_coeff_eq_trace_pow}
    \lean{MPOTensor.BNTBlockedBasisCoefficientComparison.%
        blocked_coeff_eq_ofChi_trace_pow}
    \lean{MPOTensor.BNTBlockedBasisCoefficientComparison.%
        blocked_coeff_eq_pulledBlockedChi_trace_pow}
    \leanok
    \uses{def:mpdo_bnt_blocked_basis_coefficient_comparison, def:mpdo_positive_bnt_label_chi_trace_power_form, lem:mpdo_positive_bnt_label_chi_eq_trace_matrix_pow}
    Suppose the chosen blocked-basis coefficients satisfy the comparison
    equality with the BNT structure coefficients, and the latter have a
    trace-power representation by positive diagonal matrices
    $\chi_{\alpha,\beta,\gamma}$ independent of length. Then, for every
    $n>0$ and every multiplication triple $(i,j,k)$,
    \begin{align}
        c^{(n)}_{i,j,k}
         & =
        \tr\!\left(\chi_{\sigma_n(i),\sigma_n(j),\tau_n(k)}^{\,n}\right).
        \notag
    \end{align}
\end{lemma}
\begin{proof}\leanok
    Use the comparison equality to replace the blocked-basis coefficient by the
    corresponding BNT structure coefficient, and then apply the
    trace-power formula.
    In the special case where the compared BNT coefficients are
    defined by traces of powers of $\chi$, the same conclusion
    follows directly from their definition.
\end{proof}

\begin{definition}[Diagonal matrices for blocked multiplication coefficients]%
    \label{def:mpdo_blocked_structure_chi_family}
    \lean{MPOTensor.AlgebraStructureData.BlockedStructureChiFamily}
    \leanok
    \uses{def:mpdo_blocked_structure_coefficients}
    For each positive blocking length $n$, choose a family of diagonal
    matrices indexed by triples in the disjoint union of the chosen basis
    indices of $\mathcal A_n$ and $\mathcal A_{2n}$. For a multiplication
    triple $(i,j,k)$, where $i,j$ label basis elements of $\mathcal A_n$ and
    $k$ labels a basis element of $\mathcal A_{2n}$, write its size as
    $r_{n,i,j,k}\in\N$ and its complex diagonal entries as
    $\chi_{n,i,j,k,s}$, with $0\le s<r_{n,i,j,k}$. The corresponding matrix is
    \begin{align}
        \chi_{n,i,j,k}
         & =
        \diag(\chi_{n,i,j,k,s})_{s=0}^{r_{n,i,j,k}-1}.
        \notag
    \end{align}
    Unlike the matrices indexed by fixed BNT labels in
    \cite[Theorem~4.14(ii)]{Cirac2016MPDO_arXiv}, the size and diagonal
    entries may depend on the blocking length and on the chosen basis indices.
\end{definition}
% The uniformity gap is documented in
% docs/paper-gaps/cpgsv17_blocked_chi_uniformity.tex.

\begin{lemma}[Trace of a power of a blocked diagonal matrix]%
    \label{lem:mpdo_blocked_structure_chi_trace_matrix_pow}
    \lean{MPOTensor.AlgebraStructureData.BlockedStructureChiFamily.trace_matrix_pow}
    \leanok
    \uses{def:mpdo_blocked_structure_chi_family}
    For each positive blocking length $n$, each multiplication triple $(i,j,k)$,
    and every exponent $L\in\N$,
    $\tr(\chi_{n,i,j,k}^{\,L})
        = \sum_{s=0}^{r_{n,i,j,k}-1} \chi_{n,i,j,k,s}^{\,L}$.
\end{lemma}
\begin{proof}\leanok
    Diagonal matrices raise to powers entrywise and the trace of a diagonal
    matrix is the sum of its entries.
\end{proof}

\begin{definition}[Trace-power form for blocked multiplication coefficients]%
    \label{def:mpdo_has_blocked_structure_chi_trace_power_form}
    \lean{MPOTensor.AlgebraStructureData.HasBlockedStructureChiTracePowerForm}
    \leanok
    \uses{def:mpdo_blocked_structure_coefficients, def:mpdo_blocked_structure_chi_family}
    The blocked multiplication coefficients are in trace-power form when, for
    every positive blocking length $n$ and every multiplication triple $(i,j,k)$,
    the coefficient $c^{(n)}_{i,j,k}$ equals
    $\sum_{s=0}^{r_{n,i,j,k}-1} \chi_{n,i,j,k,s}^{\,n}$.
    Equivalently, by
    Lemma~\ref{lem:mpdo_blocked_structure_chi_trace_matrix_pow}, it is the
    trace of the $n$-th power of the corresponding diagonal matrix. The exponent
    $n$ is the blocking length of the two factors in $\mathcal A_n$; the
    product is expanded in the basis of $\mathcal A_{2n}$.
\end{definition}

\begin{lemma}[Blocked coefficient as a trace power]%
    \label{lem:mpdo_blocked_structure_chi_eq_trace_matrix_pow}
    \lean{MPOTensor.AlgebraStructureData.HasBlockedStructureChiTracePowerForm.eq_trace_matrix_pow}
    \lean{MPOTensor.AlgebraStructureData.PositiveBlockedStructureChiTracePowerForm.eq_trace_pow}
    \leanok
    \uses{def:mpdo_has_blocked_structure_chi_trace_power_form, lem:mpdo_blocked_structure_chi_trace_matrix_pow}
    Under trace-power form for the blocked multiplication coefficients,
    for every positive blocking length $n$,
    $c^{(n)}_{i,j,k}=\tr(\chi_{n,i,j,k}^{\,n})$.
\end{lemma}
\begin{proof}\leanok
    Combine the defining trace-power identity with the trace formula for
    powers of a diagonal matrix.
\end{proof}

\begin{definition}[Positive diagonal matrices giving the blocked coefficients]%
    \label{def:mpdo_positive_blocked_structure_chi_trace_power_form}
    \lean{MPOTensor.AlgebraStructureData.PositiveBlockedStructureChiTracePowerForm}
    \leanok
    \uses{def:mpdo_has_blocked_structure_chi_trace_power_form, def:mpdo_blocked_structure_chi_family}
    Consider a family of blocked diagonal matrices with strictly positive
    diagonal entries satisfying
    $c^{(n)}_{i,j,k}=\tr(\chi_{n,i,j,k}^{\,n})$ for every $n>0$ and every
    multiplication triple $(i,j,k)$. This is a blocked-basis analogue of the
    positive diagonal matrices in
    \cite[Theorem~4.14(ii)]{Cirac2016MPDO_arXiv}, without requiring a single
    family indexed only by the BNT labels.
\end{definition}
% The uniformity gap is documented in
% docs/paper-gaps/cpgsv17_blocked_chi_uniformity.tex.

\begin{theorem}[Positive BNT diagonal matrices give positive blocked diagonal matrices]%
    \label{thm:mpdo_bnt_blocked_basis_to_positive_blocked_chi_trace_power_form}
    \lean{MPOTensor.BNTBlockedBasisCoefficientComparison.blockedLabel}
    \lean{MPOTensor.BNTBlockedBasisCoefficientComparison.pulledBlockedChiFamily}
    \lean{MPOTensor.BNTBlockedBasisCoefficientComparison.pulledBlockedChiFamily_toDiagonal_of_pos}
    \lean{MPOTensor.BNTBlockedBasisCoefficientComparison.%
        pulledBlockedChiFamily_toDiagonal_posEntries}
    \lean{MPOTensor.BNTBlockedBasisCoefficientComparison.pulledBlockedChi_tracePowerCoeff_of_pos}
    \lean{MPOTensor.BNTBlockedBasisCoefficientComparison.pulledBlockedChi_dim_of_pos}
    \lean{MPOTensor.BNTBlockedBasisCoefficientComparison.pulledBlockedChi_trace_matrix_pow_of_pos}
    \lean{MPOTensor.BNTBlockedBasisCoefficientComparison.%
        toPositiveBlockedStructureChiTracePowerForm}
    \leanok
    \uses{def:mpdo_bnt_blocked_basis_coefficient_comparison, def:mpdo_positive_bnt_label_chi_trace_power_form, def:mpdo_positive_blocked_structure_chi_trace_power_form}
    Suppose the chosen blocked-basis coefficients satisfy the comparison
    equality with the BNT structure coefficients, and the latter have a
    trace-power representation by positive diagonal matrices
    $\chi_{\alpha,\beta,\gamma}$ independent of length. Then there are
    positive blocked diagonal matrices giving the blocked coefficients by
    the trace-power formula.
    They are obtained by reindexing the BNT diagonal matrices along
    $\sigma_n\sqcup\tau_n$; on multiplication triples this gives
    $\chi_{n,i,j,k}=\chi_{\sigma_n(i),\sigma_n(j),\tau_n(k)}$.
\end{theorem}
\begin{proof}\leanok
    Reindex the positive diagonal matrices along the blocked-basis
    label map. Positivity is preserved by reindexing, and the trace-power
    identity follows by substituting the blocked-basis comparison into the
    BNT trace-power formula.
\end{proof}

\begin{lemma}[Constant power sums have entries zero or one]%
    \label{lem:constant_power_sum_dichotomy}
    \lean{le_one_of_forall_sum_pow_eq}
    \lean{eq_zero_or_eq_one_of_forall_sum_pow_eq}
    \lean{Complex.eq_zero_or_eq_one_of_forall_sum_pow_eq}
    \leanok
    Let $x_0,\ldots,x_{n-1}$ be non-negative real numbers whose power sums
    $\sum_k x_k^L$ take the same value at every integer exponent $L \ge 1$.
    Then every $x_k$ equals $0$ or $1$. The same holds for a finite family of
    complex numbers whose imaginary parts vanish and whose real parts are
    non-negative.
\end{lemma}
\begin{proof}\leanok
    If some entry exceeded one, its powers would grow without bound while
    all other summands stay non-negative, contradicting constancy of the
    power sums; hence every entry lies in $[0,1]$. Comparing the exponents
    one and two gives $\sum_k x_k(1-x_k) = 0$ with non-negative summands, so
    each summand vanishes.
\end{proof}

\begin{definition}[Length-independent coefficients]%
    \label{def:mpdo_bnt_label_length_independent}
    \lean{MPOTensor.BNTLabelCoefficientFamily.LengthIndependent}
    \lean{MPOTensor.BNTLabelCoefficientFamily.LengthIndependent.coeff_eq}
    \leanok
    \uses{def:mpdo_bnt_label_coefficient_family}
    The structure coefficients $c^{(L)}_{\alpha,\beta,\gamma}$ are
    \emph{length independent} when
    $c^{(L)}_{\alpha,\beta,\gamma} = c^{(1)}_{\alpha,\beta,\gamma}$ for
    every positive length $L$ and all labels; equivalently, the
    coefficients agree at any two positive lengths. This is the case
    singled out in the discussion following
    \cite[Theorem~4.14]{Cirac2016MPDO_arXiv}.
\end{definition}

\begin{theorem}[Length-independent coefficients are natural numbers]%
    \label{thm:mpdo_bnt_label_length_independent_natural}
    \lean{MPOTensor.DiagonalChiFamily.%
        entry_eq_one_of_posEntries_of_forall_tracePowerCoeff_eq}
    \lean{MPOTensor.DiagonalChiFamily.matrix_eq_one_of_forall_entry_eq_one}
    \lean{MPOTensor.BNTLabelCoefficientFamily.HasPositiveLengthChiTracePowerForm.%
        entry_eq_one_of_lengthIndependent}
    \lean{MPOTensor.BNTLabelCoefficientFamily.HasPositiveLengthChiTracePowerForm.%
        coeff_eq_dim_of_lengthIndependent}
    \lean{MPOTensor.BNTLabelCoefficientFamily.HasPositiveLengthChiTracePowerForm.%
        exists_natCast_coeff_of_lengthIndependent}
    \lean{MPOTensor.PositiveBNTLabelChiTracePowerForm.entry_eq_one_of_lengthIndependent}
    \lean{MPOTensor.PositiveBNTLabelChiTracePowerForm.coeff_eq_dim_of_lengthIndependent}
    \leanok
    \uses{def:mpdo_bnt_label_length_independent, def:mpdo_bnt_label_positive_length_chi_trace_power_form, def:mpdo_positive_bnt_label_chi_trace_power_form, lem:constant_power_sum_dichotomy}
    Suppose the BNT structure coefficients have the trace-power form
    $c^{(L)}_{\alpha,\beta,\gamma}
        = \tr(\chi_{\alpha,\beta,\gamma}^{\,L})$
    with positive diagonal matrices $\chi_{\alpha,\beta,\gamma}$, and are
    length independent. Then every diagonal entry
    $\chi_{\alpha,\beta,\gamma,r}$ lies in $\{0,1\}$, as stated in
    \cite[Section~4.5]{Cirac2016MPDO_arXiv}; by positivity every entry
    equals one. Consequently, at every positive length,
    $c^{(L)}_{\alpha,\beta,\gamma}
        = r_{\alpha,\beta,\gamma} \in \N$, where
    $r_{\alpha,\beta,\gamma}$ is the size of the diagonal matrix.
    The structure coefficient is this multiplicity, not necessarily one.
\end{theorem}
\begin{proof}\leanok
    For fixed labels, the trace-power form turns length independence into
    constancy of the power sums of the diagonal entries of
    $\chi_{\alpha,\beta,\gamma}$, so by
    Lemma~\ref{lem:constant_power_sum_dichotomy} every entry equals $0$ or
    $1$, and positivity excludes $0$. The trace of
    $\chi_{\alpha,\beta,\gamma}^{\,L}$ is then the number of diagonal
    entries.
\end{proof}

\begin{theorem}[Unit entries give length independence]%
    \label{thm:mpdo_bnt_label_length_independent_of_unit_entries}
    \lean{MPOTensor.DiagonalChiFamily.tracePowerCoeff_eq_dim_of_forall_entry_eq_one}
    \lean{MPOTensor.BNTLabelCoefficientFamily.HasPositiveLengthChiTracePowerForm.%
        lengthIndependent_of_forall_entry_eq_one}
    \leanok
    \uses{def:mpdo_bnt_label_length_independent, def:mpdo_bnt_label_positive_length_chi_trace_power_form}
    If the BNT structure coefficients have a trace-power representation by
    diagonal matrices all of whose diagonal entries equal one, then for
    every $L>0$, $c^{(L)}_{\alpha,\beta,\gamma}
        = r_{\alpha,\beta,\gamma}$, so the coefficients are length
    independent.
\end{theorem}
\begin{proof}\leanok
    The power sums of a family of ones do not depend on the exponent.
\end{proof}
